\begin{Abstract}
\selectlanguage{english}
Integrating satellite imagery into agricultural software requires consistency across geometries, dates, units and outputs. This study developed and technically assessed a modular web platform using Python, Streamlit, Google Earth Engine, Sentinel-2 and NASA POWER. Development followed a specification-driven workflow. The platform computes the Normalized Difference Vegetation Index (NDVI) and near-infrared--shortwave-infrared variants labelled Normalized Difference Water Index (NDWI) and Normalized Difference Moisture Index (NDMI). It integrates maps, time series, weather data, browser-side grid queries and exports. The documented assessment recorded 710 passing Python tests, 100\% line and branch coverage in the measured modules, three JavaScript checks and a demonstration with real services. A completed exploratory questionnaire with \UsuariosN{} responses yielded a mean rating of \UsuariosMediaEN{}/10, a median of \UsuariosMediana{} and barriers involving area selection and alert interpretation. The contribution lies in data integration, state consistency and traceable outputs. The evidence supports the demonstrated computational workflow; the questionnaire describes self-reported use, without population-level generalization. Agronomic field validation and calibration of the temporal heuristics remain future work.
\end{Abstract}
\begin{keywords}
software engineering; agricultural monitoring; remote sensing; geospatial processing; data visualization.
\end{keywords}
\selectlanguage{brazilian}

\begin{resumo}
Integrar imagens de satélite a um software agrícola exige consistência entre geometrias, datas, unidades e saídas. Este estudo desenvolveu e avaliou tecnicamente uma plataforma web modular com Python, Streamlit, Google Earth Engine, Sentinel-2 e NASA POWER. O desenvolvimento seguiu um fluxo orientado a especificações. A plataforma calcula o índice de vegetação por diferença normalizada (NDVI) e variantes infravermelho próximo--infravermelho de ondas curtas denominadas índice de água por diferença normalizada (NDWI) e índice de umidade por diferença normalizada (NDMI). Integra mapas, séries temporais, dados meteorológicos, consulta a grades no navegador e exportações. A avaliação documentada registrou 710 testes Python aprovados, com 100\% de cobertura de linhas e ramificações nos módulos medidos, três verificações JavaScript e demonstração com serviços reais. O questionário exploratório concluído, com \UsuariosN{} respostas, apresentou nota média de \UsuariosMedia{}/10, mediana \UsuariosMediana{} e barreiras de demarcação e interpretação de alertas. A contribuição está na integração de dados, na coerência do estado e na rastreabilidade das saídas. As evidências sustentam o fluxo computacional demonstrado; o questionário descreve o uso autodeclarado, sem generalização populacional. Validação agronômica em campo e calibração das heurísticas temporais permanecem como trabalhos futuros.
\end{resumo}
\begin{palavras_chave}
engenharia de software; monitoramento agrícola; sensoriamento remoto; processamento geoespacial; visualização de dados.
\end{palavras_chave}
